\newsection{Mathematische Logik und Beweisführung}{task-4}

\begin{enumerate}
  \item{
    Vera hat in diesem Wintersemester angefangen zu studieren und bereits ihre
    ersten Vorlesungen gehört.
    Sie mag dabei manche Vorlesungen lieber als andere.

    Folgende Aussagen über {\glqq}Vera mag Vorlesung ...{\grqq} sind bereits bekannt:

    {
      \renewcommand{\labelenumii}{(\roman{enumii})}
      \begin{enumerate}
        \item{Vera mag Vorlesung-1 nicht.}
        \item{Wenn Vera Vorlesung-2 nicht mag und Vorlesung-3 mag, dann mag sie auch Vorlesung-1.}
        \item{Vera mag Vorlesung-3 oder sie mag Vorlesung-2 nicht.}
      \end{enumerate}
    }

    Die Aussagen sollen als Aussagenlogische Formeln formuliert werden.

    \begin{enumerate}
      \item {
        Geben Sie geeignte Variablen und deren Bedeutung für die Aussagen an!

        \solution{
          V1: Vera mag Vorlesung-1 \\
          V2: Vera mag Vorlesung-2 \\
          V3: Vera mag Vorlesung-3
        }{\vfill}
      }
      \item {
        Formulieren Sie die Aussagen mit den Variablen als
        \textbf{syntaktisch korrekte Aussagenlogische Formeln}.
        \begin{enumerate}
          \item{
            \solution{
              $\varphi_1 = \neg V1$
            }{}
          }
          \item{
            \solution{
              $\varphi_2 = ((\neg V2 \wedge V3) \to V1)$
            }{\vspace{20pt}}
          }
          \item{
            \solution{
              $\varphi_3 = (\neg V2 \vee V3)$
            }{\vspace{20pt}}
          }
        \end{enumerate}
        \solution{}{\vspace{20pt}} % The spacing is weird, which is why we have to shift all vspaces by one item
      }
      \item {
        Finden Sie eine Bewertung, mit der alle Regeln erfüllt sind.

        \solution{
          Mit $B(V1) = 0$ und $B(V2) = B(V3) = 1$ sind alle Regeln erfüllt.
        }{\vfill}
      }
      \item {
        Ist es möglich, dass alle Regeln erfüllt sind und Vera Vorlesung-2 anders
        bewertet als Vorlesung-3?
        Begründen Sie Ihre Antwort kurz mittels der Methodik aus der Vorlesung!

        \solution{
          Nein, dies ist nicht möglich. Damit Regel (i) erfüllt ist, muss $B(V1) = 0$ sein.
          Damit dann (ii) erfüllt ist, muss die Implikatiosnannahme falsch sein.
          Dies ist der Fall, wenn mindestens entweder $B(V2) = 1$ oder $B(V3) = 0$.
          Damit allerdings gleichzeitig (iii) erfüllt ist, muss die Bewertung der
          Variable, die in (ii) dafür sorgt, dass die Implikationsannahme falsch ist,
          die Disjunktion in (iii) erfüllen, also muss $B(V2) = 0$ oder $B(V3) = 1$ sein.

          Die folgende Wahrheitstabelle zeigt den Zusammenhang anschaulich:

          \vspace{10pt}
          \setlength{\tabcolsep}{5pt}
          \begin{center}
            \begin{tabular}{c|c||c|c||c|c}
              $V2$ & $V3$ & $(\neg V2 \wedge V3)$ & $(\neg V2 \vee V3)$ & $\varphi_2$ & $\varphi_3$ \\ \hline\hline
              0 & 0 & 0 & 1 & 1 & 1 \\
              0 & 1 & 1 & 1 & 0 & 1 \\
              1 & 0 & 0 & 0 & 1 & 0 \\
              1 & 1 & 0 & 1 & 1 & 1
            \end{tabular}
          \end{center}

          Wie zu sehen ist, ermöglicht nur eine gleiche Bewertung von $V2$ und $V3$,
          dass sowohol $\varphi_2$ als auch $\varphi_3$ erfüllt sind.
        }{\vfill}
      }
    \end{enumerate}
  }
  \newpage
  \item{
    Nun betrachten wir den korrekten Umgang mit Wertetabellen zur Aussagenlogik!
    Folgende Aussagenlogische Formeln sind gegeben:

    \[
      \varphi_1 = (X \wedge (Y \vee Z)) \quad
      \varphi_2 = ((X \wedge Y) \vee (X \wedge Z)) \quad
      \varphi = (\varphi_1 \leftrightarrow \varphi_2)
    \]

    \vspace{10pt}
    \begin{enumerate}
      \item {
        Welche allgemeingültige Aussagen wird durch $(\varphi_1 \leftrightarrow \varphi_2)$
        ausgedrückt?

        \solution{
          Die Allgemeingültigkeit von $(\varphi_1 \leftrightarrow \varphi_2)$
          drückt das Distributivgesetz der Disjunktion aus.
        }{\vspace{40pt}}
      }
      \item {
        Werten Sie die folgende Wahrheitstabelle aus, indem Sie die leeren
        Spalten mit \textbf{syntaktisch korrekten} Hilfsaussagen füllen!

        \vspace{10pt}
        \setlength{\tabcolsep}{10pt}
        \renewcommand{\arraystretch}{1.5}
        \begin{center}
          \begin{tabular}{c|c|c||c|c|c||c|c|c}
            $X$ & $Y$ & $Z$ & $(Y \vee Z)$ & \solution{$(X \wedge Y)$}{\hspace{60pt}} & \solution{$(X \wedge Z)$}{\hspace{60pt}} & $\varphi_1$ & $\varphi_2$ & $\varphi$ \\ \hline\hline
            0 & 0 & 0 & \solution{0}{} & \solution{0}{} & \solution{0}{} &\solution{0}{}  & \solution{0}{} & \solution{1}{} \\ \hline
            0 & 0 & 1 & \solution{1}{} & \solution{0}{} & \solution{0}{} &\solution{0}{}  & \solution{0}{} & \solution{1}{} \\ \hline
            0 & 1 & 0 & \solution{1}{} & \solution{0}{} & \solution{0}{} &\solution{0}{}  & \solution{0}{} & \solution{1}{} \\ \hline
            0 & 1 & 1 & \solution{1}{} & \solution{0}{} & \solution{0}{} &\solution{0}{}  & \solution{0}{} & \solution{1}{} \\ \hline
            1 & 0 & 0 & \solution{0}{} & \solution{0}{} & \solution{0}{} &\solution{0}{}  & \solution{0}{} & \solution{1}{} \\ \hline
            1 & 0 & 1 & \solution{1}{} & \solution{0}{} & \solution{1}{} &\solution{1}{}  & \solution{1}{} & \solution{1}{} \\ \hline
            1 & 1 & 0 & \solution{1}{} & \solution{1}{} & \solution{0}{} &\solution{1}{}  & \solution{1}{} & \solution{1}{} \\ \hline
            1 & 1 & 1 & \solution{1}{} & \solution{1}{} & \solution{1}{} &\solution{1}{}  & \solution{1}{} & \solution{1}{} \\ \hline
          \end{tabular}
        \end{center}
        \vspace{20pt}
      }
      \item {
        Wie lässt sich anhand dieser Wertetabelle erkennen, dass $\varphi$
        eine Tautologie ist?
        Begründen Sie kurz.

        \solution{
          Es lässt sich anhand der Wertetabelle erkennen, dass $\varphi$ immer
          wahr (1) und somit unabhängig von den Bewertungen von $X, Y$ und $Z$
          ist.
          Das macht $\varphi$ definitionsgemäß zu einer Tautologie.
        }{}
      }
    \end{enumerate}
  }
  \newpage
  \item {
    Wir betrachten das Resolutionskalkül für Aussgnelogische Formeln und
    dessen Ableitungsregeln!

    Doktor Halbschlau hat sich für das Resolutionskalkül der Aussagenlogik
    folgende neue Ableitungsregeln ausgedacht.

    Aus zwei beliebige Klauseln $C_1, C_2$ entsteht immer und ohne Bedingungen
    in einem einzigen Resolutionsschritt eine neue Klausel
    \[
      C = C_1 \cup C_2
    \]

    Die Abbildung zeigt ein ganz einfaches Beispiel der Anwendung dieser Regel:

    \begin{center}
      \begin{tikzpicture}[distance = 2cm]
        \node (c1) {$C_1 = \{\neg X, Z\}$};
        \node (c) [right=of c1, below=1.5cm] {$C = \{\neg X, Z, Y, W\}$};
        \node (c2) [right=of c, above=1.5cm] {$C_2 = \{\neg X, Y, W\}$};

        \draw (c1) -- (c);
        \draw (c2) -- (c);
        % below=2cm
      \end{tikzpicture}
    \end{center}

    \vspace{20pt}
    Seien $C_1 = \{\neg x_1, x_2\},\: C_2 = \{x_1\}$ und $C_3 = \{\neg x_1, x_3\}$ Klauseln.

    \begin{enumerate}
      \item {
        Für eine Klauselmenge $K$ ist Res$^1(K)$ wie folgt definiert.
        Geben Sie konkret die neuen Resolventen an, die sich aus dem
        Resolutionskalkül von Doktor Halbschlau für $C_1, C_2$ und $C_3$
        ableiten lassen!

        \vspace{10pt}
        \[
          \solution{}{\hspace{-120pt}} \text{Res}^1(K) \coloneq \text{Res}(K)
            = K \cup \{ \solution{\{\neg x_1, x_1\}, \{\neg x_1, x_3\}, \{x_1, x_3\} \} }{}
        \]
        \vspace{10pt}
      }
      \item {
        Geben Sie nun auch die Definition von Res$^2(K)$ \textit{und}
        die konkreten Resolventen, die sich im 2. Schritt ableiten lassen, an.

        \vspace{10pt}
        \[
          \solution{}{\hspace{-60pt}} \text{Res}^2(K) \coloneq \solution{\text{Res}(\text{Res}^1(K))}{\hspace{120pt}} = \solution{\text{Res}^1(K) \cup \{\}}{\hspace{60pt} \cup \; \{}
        \]
        \vspace{10pt}
      }
      \item {
        Ab welchem $n \in \mathbb{N}$ gilt Res$^n(K) =$ Res$^*(K)$?
        Geben Sie jenes $n$ an und begründen Sie in einem \textbf{einzigen Satz}!

        \solution{
          Ab $n=1 \in \mathbb{N}$ gilt Res$^n(K) =$ Res$^*(K)$, da im 2. Schritt (s.o.)
          keine neuen Resolventen ableitbar waren, also Res$^1(K) =$ Res$^2(K)$.
        }{}
      }
    \end{enumerate}
  }
  \newpage
  \item {
    Wir wollen uns mit der Korrektheit und der Vollständigkeit des
    vorhin vorgestellten Resolutionskalküls von Doktor Halbschlau beschäftigen.

    \begin{enumerate}
      \item {
        Beweisen Sie allgemein anhand einer Bewertung $B$ und deren Bedeutung,
        dass, wenn die Klauseln $\{C_1, C_2\}$ erfüllt sind, auch immer die
        Klauselmenge $\{C_1, C_2, C\}$ erfüllt ist.

        \solution{
          Wenn die Klauselmenge $\{C_1, C_2\}$ erfüllt ist, müssen definitionsgemäß
          jeweils alle Klauseln, also $C_1$ und $C_2$ erfüllt sein.
          Dafür muss mindestens ein Literal aus je einer Klausel erfüllt sein.
          Damit dies der Fall ist, muss für eine Bewertung eines Literals $l \in C$
          $B(l) = 1$, sofern $l$ ein positives Literal ist, oder $B(l) = 0$,
          sofern $l$ ein negatives Literal ist, gelten.

          Die abgeleitete Klausel $C$ enthält nach der Ableitungsregel des
          Resolutionskalküls alle Literale aus $C_1$ und $C_2$.
          Da bereits die Klauselmenge $\{C_1, C_2\}$ erfüllt ist und somit jeweils
          mindestens ein Literal aus $C_1$ und $C_2$ erfüllt sind, ist auch
          mindestens ein Literal aus $C$ erfüllt.
          Somit ist $C$ und damit auch die Klauselmenge $\{C_1, C_2, C\}$ erfüllt.
        }{\vfill}
      }
      \item {
        Beweisen oder widerlegen Sie in einem \textbf{einzigen Satz}, dass
        die Ableitungsregeln des Resolutionskalküls \textit{korrekt} im Sinne der
        Vorlesung sind.

        \solution{
          Da, wie oben gezeigt, alle abgeleiteten Klauseln aus zwei erfüllten Klauseln
          ebenfalls erfüllt sind, werden nur wahre Aussagen abgeleitet und
          das Resolutionskalkül ist definitionsgemäß korrekt.
        }{\vspace{50pt}}
      }
      \item {
        Zeigen Sie anhand eines kleinen Beispiels, dass das Resolutionskalkül
        nicht \textit{vollständig} im Sinne der Vorlesung ist.

        \solution{
          Seien $C_1 = \{\neg x_1, x_2\}$ und $C_1 = \{x_1, x_2\}$ Klauseln.

          Mittels dem aus der Vorlesung bekannten Resolutionskalkül lässt sich
          die wahre Klausel $C = \{x_2\}$ ableiten.
          Aus dem Resolutionskalkül von Doktor Halbschlau lässt sich allerdings
          nur die Klausel $C' = \{\neg x_1, x_1, x_2\}$ ableiten.

          Somit lassen sich aus dem gegebenen Resolutionskalkül nicht alle
          wahren Aussagen ableiten und das Resolutionskalkül ist nicht vollständig
          im Sinne der Vorlesung.
        }{\vfill \vfill}
      }
    \end{enumerate}
  }
\end{enumerate}